% ============================================================================
% Fichier  : partie5/P05_C15_rapport_temoignage.tex
% Titre    : Le Rapport et le Témoignage
% Auteur   : MINKA MI NGUIDJOÏ Thierry Emmanuel
% Version  : 1.0 -- novembre 2026
% Corrigés : corriges/Corr_C15.tex
% ============================================================================
% Sources : ISO/IEC 27042 (interprétation et rapport, paraphrasée) ; ISO/IEC
% 17025 (compétence des laboratoires) ; loi n° 2010/012, art. 56 (rappel).
% Faits nouveaux (bible) : rapport d'expertise de M. ABOUBAKAR déposé le
% 28/05/2025, 46 pages et 9 annexes ; audience du 15/02/2026 devant le
% tribunal de grande instance de la Mifi (Bafoussam) : deux titulaires de
% portefeuilles poursuivis pour recel ; l'expert n'attribue pas l'intrusion.
% ============================================================================

\chapter{Le Rapport et le Témoignage}
\label{chap:rapport}
\chaptermark{Le rapport et le témoignage}

\epigraph{%
    « Un rapport que seul son auteur comprend\\
    n'a convaincu personne. »%
}{Principe du cours}

\niveauChapitre{Intermédiaire}{%
    Maîtrise du Chapitre~\ref{chap:enquete}
}

\begin{boxobjectifs}
\begin{itemize}
    \item Adapter un rapport à son destinataire~: direction, autorité
          judiciaire, équipe technique.
    \item Structurer un rapport d'expertise de façon qu'il soit vérifiable.
    \item Écrire chaque conclusion avec sa constatation, son analyse et son
          niveau de confiance.
    \item Dire les limites de l'examen sans les cacher ni s'y réfugier.
    \item Préparer une déposition et répondre à une question hostile.
\end{itemize}
\end{boxobjectifs}

\bigskip

% ============================================================================
\section*{Avant de commencer : une question simple}
% ============================================================================

Bafoussam, tribunal de grande instance de la Mifi, 15 février 2026. Deux
titulaires de portefeuilles de paiement mobile, qui ont reçu une partie des
fonds détournés, sont poursuivis pour recel. M.~ABOUBAKAR, l'expert requis
par le procureur, dépose à la barre. L'avocat de la défense lui pose une
question simple~: « Monsieur l'expert, pouvez-vous dire au tribunal que mes
clients ont piraté NYANGA CRÉDIT~? » La réponse tient en une phrase, et
l'expert l'a préparée.

\begin{boxalert}
\textbf{Violation typique.} Laisser un rapport, ou une déposition, aller
au-delà de ce que les pièces établissent, parce que le destinataire attend
une réponse ferme.

\smallskip
\textbf{Règle~:} l'expert répond à la question posée, avec ce que les pièces
établissent, et dit clairement ce qu'elles n'établissent pas. Une réponse
\NEK{} bien formulée sert la justice mieux qu'une certitude fragile.
\end{boxalert}

\separateur

% ============================================================================
\section{Trois Destinataires, Trois Rapports}
\label{sec:c15-destinataires}
% ============================================================================

\begin{tableaucours}{L{2.6cm}L{3.4cm}L{3.2cm}L{3.0cm}}{%
    \tableauHeader{Destinataire} &
    \tableauHeader{Ce qu'il doit décider} &
    \tableauHeader{Forme} &
    \tableauHeader{Écueil}
}
Direction & Fermer la porte, financer, informer & Note de deux pages,
recommandations datées & Trop de technique \\
Autorité judiciaire & Établir des faits, qualifier & Rapport d'expertise
complet, annexes & Sortir de la mission \\
Équipe technique & Corriger, détecter & Rapport détaillé, indicateurs,
règles & Pas de priorités \\
\end{tableaucours}
\vspace{-0.8em}
\begin{center}{\small\itshape Tableau~15.1. Trois rapports pour une même enquête}\end{center}

Les trois rapports reposent sur les mêmes constatations. Ils diffèrent par
ce qu'ils retiennent et par leur langue. La note à la direction ne contient
pas de vecteur CVSS~; le rapport d'expertise ne contient pas de
recommandation budgétaire~; le rapport technique contient les deux.

% ============================================================================
\section{Le Rapport d'Expertise}
\label{sec:c15-expertise}
% ============================================================================

L'ISO/IEC~27042~\cite{iso27042_2015} demande que le rapport permette à un
tiers compétent de comprendre ce qui a été fait, d'en évaluer la validité et
de le refaire (paraphrase). Le plan suivant répond à cette exigence.

\begin{tableaucours}{cL{3.6cm}L{8.0cm}}{%
    \tableauHeader{} &
    \tableauHeader{Partie} &
    \tableauHeader{Contenu}
}
1 & Mission & Autorité requérante, fondement (art.~56 de la loi
n°~2010/012), questions posées, mot pour mot \\
2 & Pièces reçues & Scellés, état à réception, empreintes vérifiées \\
3 & Opérations & Chaque opération, datée, avec outils, versions et copies
utilisées \\
4 & Constatations & Les faits, sans interprétation, avec leur source \\
5 & Discussion & Les interprétations, les autres explications examinées,
le niveau de confiance \\
6 & Limites & Ce qui n'a pas pu être examiné, et pourquoi \\
7 & Réponses aux questions & Une réponse par question, cotée \\
8 & Annexes & Sorties brutes, tableaux, chronologie, formulaires \\
\end{tableaucours}
\vspace{-0.8em}
\begin{center}{\small\itshape Tableau~15.2. Le plan d'un rapport d'expertise}\end{center}

\subsection{Écrire une conclusion}

Chaque conclusion du rapport suit le même schéma~: la constatation, avec sa
source~; l'analyse, avec les autres explications examinées~; le niveau de
confiance, et ce qui le fonde. Le vocabulaire de la confiance doit être
défini dans le rapport et tenu partout.

\begin{tableaucours}{L{2.8cm}L{9.4cm}}{%
    \tableauHeader{Niveau} &
    \tableauHeader{Usage dans ce livre}
}
Élevée & Plusieurs pièces indépendantes concordent, aucune explication
concurrente plausible ne subsiste \\
Moyenne & Les pièces concordent mais l'une est fragile, ou une explication
concurrente reste possible \\
Faible & Une seule pièce, ou des explications concurrentes aussi plausibles \\
\end{tableaucours}
\vspace{-0.8em}
\begin{center}{\small\itshape Tableau~15.3. Les niveaux de confiance}\end{center}

\begin{boxscenario}{Opération NYANGA. Une réponse du rapport d'expertise}
\textit{Question~3 de la mission~: « Dire par quel moyen l'accès au système
d'information de NYANGA CRÉDIT a été obtenu. »}

\medskip
\textbf{Constatation.} La passerelle d'accès à distance fonctionnait le
2~mars 2025 avec le micrologiciel 7.2.1 (scellé NC-02, export de
configuration). Le bulletin de l'éditeur du 12~décembre 2024 désigne cette
version comme vulnérable à un contournement de l'authentification. Le
journal de la passerelle (NC-02) enregistre, le 2~mars entre 23~h~46 et
23~h~48, quatre requêtes depuis l'adresse 203.0.113.18 conformes aux
indicateurs publiés par l'éditeur, suivies de l'ouverture d'une session du
compte \texttt{svc\_backup} sans authentification préalable.

\textbf{Discussion.} Deux autres explications ont été examinées~:
l'utilisation d'un mot de passe connu, qui n'explique pas la session sans
authentification~; le vol d'une session existante, écarté faute de session
légitime dans les neuf jours précédents.

\textbf{Réponse} (confiance élevée). L'accès a été obtenu par l'exploitation
d'une vulnérabilité de la passerelle permettant de contourner
l'authentification.
\end{boxscenario}

\subsection{Dire les limites}

La partie consacrée aux limites protège l'expert autant que le lecteur.
Elle énumère ce que l'examen n'a pas couvert~: la mémoire de SRV-CPT01,
jamais capturée~; le système de la passerelle, illisible sans l'éditeur~;
la période du 3 au 8~mars, sans trace~; la provenance des journaux de
l'interface de paiement. Pour chaque limite, elle dit l'effet sur les
conclusions. Une limite non écrite devient, à l'audience, une faille de
l'expertise.

% ============================================================================
\section{L'Indépendance et la Compétence}
\label{sec:c15-independance}
% ============================================================================

L'expert requis tient sa légitimité de deux qualités~: son indépendance à
l'égard des parties, et sa compétence. La première exclut d'avoir travaillé
pour l'une d'elles sur l'affaire. C'est pourquoi le procureur n'a pas requis
Mme~MBARGA, mandatée et payée par NYANGA CRÉDIT, mais M.~ABOUBAKAR. La
seconde se démontre par la formation, l'expérience et la méthode~; pour un
laboratoire, l'accréditation selon l'ISO/IEC~17025~\cite{iso17025_2017}
atteste de la compétence technique et de la maîtrise des méthodes.

\begin{boxinfo}
\textbf{L'enquêteur privé comme témoin.} Mme~MBARGA peut être entendue comme
témoin des opérations qu'elle a réalisées~: acquisitions, chaîne de
possession, premières constatations. Elle ne se substitue pas à l'expert,
et l'expert ne se contente pas de reprendre ses conclusions~: il refait ou
contrôle ce qui fonde les siennes.
\end{boxinfo}

% ============================================================================
\section{Déposer à l'Audience}
\label{sec:c15-audience}
% ============================================================================

L'audience n'est pas un séminaire technique. Le tribunal, les avocats et
les prévenus n'ont pas lu les annexes. L'expert y gagne à préparer trois
choses~: une chronologie d'une page, lisible à distance~; une définition en
une phrase de chaque terme technique qu'il emploiera~; et la liste des
questions hostiles probables, avec leur réponse.

\begin{tableaucours}{L{5.4cm}L{6.8cm}}{%
    \tableauHeader{Question hostile} &
    \tableauHeader{Réponse fondée sur le dossier}
}
« Vos outils ne sont-ils pas faillibles~? » & Les conclusions principales
reposent sur deux sources indépendantes, lues avec des outils dont les
versions et les sorties sont en annexe \\
« Le disque est resté une nuit sur un bureau. » & L'intégrité est établie
par les empreintes~; la conclusion qu'il soutient est confirmée par une
autre pièce, à la chaîne complète \\
« Pouvez-vous exclure une erreur d'horloge~? » & Les dérives ont été
mesurées et corrigées~; les écarts restants sont inférieurs à la minute \\
« Mes clients ont-ils piraté l'établissement~? » & Les pièces ne permettent
pas de le dire (voir ci-dessous) \\
\end{tableaucours}
\vspace{-0.8em}
\begin{center}{\small\itshape Tableau~15.4. Préparer les questions hostiles}\end{center}

\begin{boxscenario}{Opération NYANGA. La réponse de l'expert}
\textbf{Réponse de M.~ABOUBAKAR.} « Non, Maître. Les pièces que j'ai
examinées établissent comment le système a été pénétré et comment les fonds
ont été transférés. Elles ne permettent pas d'identifier les personnes qui
ont commis l'intrusion. Elles établissent, par les relevés de l'opérateur,
que deux des onze portefeuilles qui ont reçu les fonds sont ouverts au nom
de vos clients. Ce que vos clients savaient de l'origine de ces fonds n'est
pas une question technique. »

\medskip
\textbf{Analyse.} L'expert répond à la question, sépare ce qui est établi de
ce qui ne l'est pas, et laisse au tribunal ce qui relève du tribunal.
L'avocat n'obtient ni une attribution qu'il pourrait contester, ni un refus
de répondre qu'il pourrait exploiter.
\end{boxscenario}

\begin{boxlecture}{Les réponses de l'expertise}
\begin{tableaucours}{L{7.6cm}c}{%
    \tableauHeader{Question de la mission} &
    \tableauHeader{Réponse}
}
Moyen de l'accès & \OUI{} \\
Séquence des opérations de l'attaquant & \OUI{} (sauf 3 au 8~mars, \NEK{}) \\
Données copiées & \OUI{} \\
Destination des fonds & \OUI{} (onze portefeuilles) \\
Identité des auteurs de l'intrusion & \NEK{} \\
\end{tableaucours}

\medskip
\textbf{Lecture.} Quatre réponses établies, une non tranchée. Le \NEK{} sur
les auteurs n'affaiblit pas l'expertise~; il en marque l'honnêteté. Les
poursuites pour recel reposent sur d'autres éléments, que la justice réunit
par d'autres moyens.
\end{boxlecture}

% ============================================================================
\section{Ce que le Rapport demande à chaque Métier}
\label{sec:c15-metiers}
% ============================================================================

\begin{boxmetier}{Application OPJ/Expert Judiciaire}
\textbf{Rédiger et défendre une expertise}
\begin{enumerate}
    \item Recopier les questions de la mission, mot pour mot.
    \item Séparer, pour chaque conclusion, constatation, discussion et
          réponse.
    \item Définir les niveaux de confiance et les tenir.
    \item Écrire les limites et leur effet.
    \item Préparer la chronologie d'une page et les questions hostiles.
\end{enumerate}
\end{boxmetier}

\begin{boxmetier}{Analyste en réponse à incident}
Écrire le rapport technique pour qu'il puisse devenir une pièce de
procédure~: chaque affirmation sourcée, chaque opération datée. Il le
deviendra plus souvent qu'on ne le prévoit.
\end{boxmetier}

\begin{boxmetier}{RSSI}
Lire les limites du rapport avant ses conclusions~: elles disent ce que
l'organisation doit encore préparer pour la prochaine enquête.
\end{boxmetier}

\separateur

\begin{boxresume}
\textbf{Ce qu'il faut retenir de ce chapitre~:}
\begin{itemize}
    \item Une même enquête donne trois rapports, pour la direction, la
          justice et l'équipe technique.
    \item Le rapport d'expertise recopie la mission, décrit les opérations,
          sépare constatations et discussion, écrit ses limites et répond
          question par question.
    \item Les niveaux de confiance se définissent et se tiennent.
    \item L'indépendance exclut d'avoir travaillé pour une partie~; la
          compétence se démontre.
    \item À l'audience, l'expert répond à la question posée et laisse au
          tribunal ce qui lui revient.
\end{itemize}
\end{boxresume}

\begin{boxinfo}
\textbf{Points de contrôle.} Vous avez compris ce chapitre si vous savez~:
\begin{itemize}
    \item rédiger une réponse d'expertise en trois temps~;
    \item écrire une section de limites et dire l'effet de chacune~;
    \item répondre à une question hostile sans dépasser les pièces.
\end{itemize}
\end{boxinfo}

% ============================================================================
\section*{Exercices}
% ============================================================================

\begin{boxexo}[title={\ding{45}\quad Exercice~15.1. Quel rapport~? \niveaubase}]
Dites dans quel rapport (direction, justice, technique) figure chacun des
éléments suivants~: (a) le vecteur CVSS de VULN-2024-4471~; (b) le coût de
l'incident~; (c) la sortie brute de \texttt{windows.netscan}~; (d) la
recommandation sur la conservation des journaux~; (e) les questions de la
mission.
\end{boxexo}

\begin{boxexo}[title={\ding{45}\quad Exercice~15.2. Réécrire une conclusion \niveaumoyen}]
Réécrivez en trois temps (constatation, discussion, réponse cotée) la phrase
suivante, tirée d'un brouillon de rapport~: « Il est évident que les
pirates ont volé le mot de passe de M.~KAMGA pour vider les comptes. »
\end{boxexo}

\begin{boxexo}[title={\ding{45}\quad Exercice~15.3. La section des limites \niveaumoyen}]
Rédigez la section « Limites » du rapport d'expertise de NYANGA. Pour chaque
limite, indiquez la ou les conclusions affectées et la manière dont elles le
sont.
\end{boxexo}

\begin{boxexo}[title={\ding{45}\quad Exercice~15.4. Trois questions hostiles \niveauexpert}]
Préparez la réponse de l'expert aux trois questions suivantes, en vous
appuyant sur le dossier~: (a) « Pourquoi n'avez-vous pas examiné la
passerelle elle-même~? »~; (b) « Les journaux de l'interface de paiement
vous ont été transmis par courriel par une société privée~: comment savoir
qu'ils sont authentiques~? »~; (c) « Votre rapport reprend-il simplement
celui du cabinet payé par la victime~? »
\end{boxexo}

\bigskip

\begin{boxinfo}
\textbf{Pour aller plus loin.}
\begin{itemize}
    \item L'ISO/IEC~27042~\cite{iso27042_2015}, sur l'interprétation et la
          présentation des résultats.
    \item L'ISO/IEC~17025~\cite{iso17025_2017}, pour la compétence des
          laboratoires.
\end{itemize}
\end{boxinfo}

\begin{boxinfo}
\textbf{Transition.} Le rapport dit ce que l'enquête sait et ce qu'elle ne
sait pas. Peut-on mesurer cette part d'ignorance, la comparer d'une enquête
à l'autre, la réduire méthodiquement~? C'est la question de recherche du
Chapitre~\ref{chap:auditologie}.
\end{boxinfo}
